\documentclass[aspectratio=169]{beamer}
%--- Configuración del Tema y Colores ---
\usetheme{Madrid}
\usecolortheme{whale}
\setbeamertemplate{navigation symbols}{}
\setbeamertemplate{caption}[numbered]

%--- Paquetes ---
\usepackage[utf8]{inputenc}
\usepackage[spanish]{babel}
\usepackage{graphicx}
\usepackage{booktabs}
\usepackage{tikz}
\usepackage{fontawesome5}
\usepackage{siunitx}
\usetikzlibrary{shapes.geometric, arrows, positioning, calc, patterns}

%--- Metadatos del Documento ---
\title[PM2]{Programación Matemática 2}
\subtitle{Clase 8: Hardware para IA - CPU vs GPU vs TPU vs LPU}
\author[Luis Alfredo Alvarado]{Prof. Luis Alfredo Alvarado Rodríguez}
\institute[ECFM - USAC]{
    Escuela de Ciencias Físicas y Matemáticas \\
    Universidad de San Carlos de Guatemala
}
\date{Primer Semestre, 2026}

\begin{document}

%--- Diapositiva de Título ---
\begin{frame}
    \titlepage
\end{frame}

%--- Diapositiva: Tabla de Contenidos ---
\begin{frame}{Agenda de Hoy}
    \tableofcontents
\end{frame}

%===========================================================
% SECCIÓN 1: INTRODUCCIÓN
%===========================================================
\section{¿Por Qué Importa el Hardware en IA?}

\begin{frame}{El Cuello de Botella Computacional}
    \begin{columns}
        \column{0.55\textwidth}
        \textbf{La IA moderna es intensiva en cómputo:}
        \begin{itemize}
            \item GPT-4: estimado en \textbf{1.8 trillones} de parámetros
            \item Entrenar LLaMA-2 70B: \textbf{1.7 millones} de horas GPU
            \item Una inferencia de imagen con Stable Diffusion: \textbf{billones} de operaciones
        \end{itemize}
        
        \vspace{0.3cm}
        \textbf{El hardware correcto puede significar:}
        \begin{itemize}
            \item Días vs meses de entrenamiento
            \item Milisegundos vs segundos de inferencia
            \item \$100 vs \$100,000 en costos
        \end{itemize}
        
        \column{0.45\textwidth}
        \centering
        \begin{tikzpicture}
            \draw[thick, ->] (0,0) -- (4.5,0) node[right, font=\scriptsize] {Año};
            \draw[thick, ->] (0,0) -- (0,3.5) node[above, font=\scriptsize] {Cómputo};
            
            \draw[blue!70, ultra thick] (0.5,0.2) .. controls (2,0.5) and (3,1.5) .. (4,3.2);
            
            \node[font=\tiny] at (1,0.3) {2012};
            \node[font=\tiny] at (2.5,1) {2018};
            \node[font=\tiny] at (4,3.5) {2024};
            
            \node[font=\scriptsize, text width=2cm, align=center] at (2.5,2.5) {Crecimiento exponencial};
        \end{tikzpicture}
    \end{columns}
\end{frame}

\begin{frame}{Los Cuatro Protagonistas}
    \centering
    \begin{tikzpicture}[node distance=1.5cm]
        \tikzstyle{chip} = [rectangle, rounded corners, minimum width=2.8cm, minimum height=1.5cm, text centered, draw=black, font=\small\bfseries]
        
        \node (cpu) [chip, fill=blue!25] {CPU\\{\scriptsize Central Processing Unit}};
        \node (gpu) [chip, right=1cm of cpu, fill=green!25] {GPU\\{\scriptsize Graphics Processing Unit}};
        \node (tpu) [chip, right=1cm of gpu, fill=orange!25] {TPU\\{\scriptsize Tensor Processing Unit}};
        \node (lpu) [chip, right=1cm of tpu, fill=purple!25] {LPU\\{\scriptsize Language Processing Unit}};
        
        \node[below=0.3cm of cpu, font=\tiny, text width=2.5cm, align=center] {Intel, AMD\\Propósito general};
        \node[below=0.3cm of gpu, font=\tiny, text width=2.5cm, align=center] {NVIDIA, AMD\\Paralelismo masivo};
        \node[below=0.3cm of tpu, font=\tiny, text width=2.5cm, align=center] {Google\\Tensores/Matrices};
        \node[below=0.3cm of lpu, font=\tiny, text width=2.5cm, align=center] {Groq\\Inferencia LLM};
    \end{tikzpicture}
    
    \vspace{0.5cm}
    \begin{block}{Pregunta Clave}
        ¿Por qué no usar CPUs para todo? La respuesta está en la \textbf{arquitectura} y el \textbf{tipo de operaciones} que cada chip optimiza.
    \end{block}
\end{frame}

%===========================================================
% SECCIÓN 2: CPU
%===========================================================
\section{CPU: El Generalista}

\begin{frame}{Arquitectura de una CPU}
    \begin{columns}
        \column{0.5\textwidth}
        \centering
        \begin{tikzpicture}[scale=0.8]
            \tikzstyle{core} = [rectangle, minimum width=1.8cm, minimum height=1.8cm, draw=black, fill=blue!20]
            \tikzstyle{cache} = [rectangle, minimum width=4cm, minimum height=0.6cm, draw=black, fill=yellow!20]
            
            % Cores
            \node (c1) [core] at (0,2) {Core 1};
            \node (c2) [core] at (2.2,2) {Core 2};
            \node (c3) [core] at (0,0) {Core 3};
            \node (c4) [core] at (2.2,0) {Core 4};
            
            % Cache
            \node (l3) [cache] at (1.1,-1.5) {L3 Cache (compartida)};
            
            % Memory controller
            \node (mc) [rectangle, minimum width=4cm, minimum height=0.5cm, draw=black, fill=gray!20] at (1.1,-2.3) {\scriptsize Memory Controller};
            
            \draw[thick] (c1.south) -- ++(0,-0.3);
            \draw[thick] (c2.south) -- ++(0,-0.3);
            \draw[thick] (c3.south) -- (c3.south |- l3.north);
            \draw[thick] (c4.south) -- (c4.south |- l3.north);
        \end{tikzpicture}
        
        \column{0.5\textwidth}
        \textbf{Características:}
        \begin{itemize}
            \item Pocos núcleos potentes (4-128)
            \item Alta frecuencia (3-5 GHz)
            \item Gran caché por núcleo
            \item Ejecución fuera de orden
            \item Predicción de saltos sofisticada
            \item Optimizado para \textbf{latencia}
        \end{itemize}
        
        \vspace{0.3cm}
        \textbf{Filosofía:} \\
        ``Hacer una tarea muy rápido''
    \end{columns}
\end{frame}

\begin{frame}{CPU: Fortalezas y Limitaciones para IA}
    \begin{columns}[t]
        \column{0.48\textwidth}
        \begin{block}{Fortalezas \faThumbsUp}
            \begin{itemize}
                \item Excelente para código secuencial
                \item Versatilidad total
                \item Disponibilidad universal
                \item Ideal para preprocesamiento
            \end{itemize}
        \end{block}
        
        \column{0.48\textwidth}
        \begin{alertblock}{Limitaciones \faThumbsDown}
            \begin{itemize}
                \item Pocos núcleos = bajo paralelismo
                \item Ineficiente para matrices grandes
                \item No escala bien para deep learning
            \end{itemize}
        \end{alertblock}
    \end{columns}
    
    \vspace{0.3cm}
    \begin{exampleblock}{Cuándo usar CPU para IA}
        Inferencia de modelos pequeños, preprocesamiento de datos, lógica de control, y cuando no hay GPU disponible.
    \end{exampleblock}
\end{frame}

%===========================================================
% SECCIÓN 3: GPU
%===========================================================
\section{GPU: El Motor del Deep Learning}

\begin{frame}{De Gráficos a Inteligencia Artificial}
    \begin{columns}
        \column{0.5\textwidth}
        \textbf{Historia:}
        \begin{itemize}
            \item Diseñadas para renderizar gráficos
            \item 2012: AlexNet usa GPUs → revolución
            \item CUDA (2006): programación general en GPU
            \item Hoy: arquitecturas específicas para IA
        \end{itemize}
        
        \vspace{0.3cm}
        \textbf{¿Por qué funcionan para IA?}
        \begin{itemize}
            \item Renderizar = multiplicar matrices
            \item Deep Learning = multiplicar matrices
            \item ¡La misma operación fundamental!
        \end{itemize}
        
        \column{0.5\textwidth}
        \centering
        \begin{tikzpicture}
            % Pixel grid
            \draw[step=0.3cm, gray, very thin] (0,0) grid (2.4,2.4);
            \node[below, font=\scriptsize] at (1.2,-0.2) {Píxeles (paralelos)};
            
            % Arrow
            \draw[ultra thick, ->] (2.8,1.2) -- (3.5,1.2);
            
            % Neural network representation
            \foreach \y in {0.3, 0.9, 1.5, 2.1} {
                \draw[fill=green!30] (3.8,\y) circle (0.15);
            }
            \foreach \y in {0.6, 1.2, 1.8} {
                \draw[fill=blue!30] (4.6,\y) circle (0.15);
            }
            \foreach \y in {0.3, 0.9, 1.5, 2.1} {
                \foreach \yy in {0.6, 1.2, 1.8} {
                    \draw[gray, thin] (3.95,\y) -- (4.45,\yy);
                }
            }
            \node[below, font=\scriptsize] at (4.2,-0.2) {Neuronas (paralelas)};
        \end{tikzpicture}
    \end{columns}
\end{frame}

\begin{frame}{Arquitectura de una GPU}
    \centering
    \begin{tikzpicture}[scale=0.75]
        \tikzstyle{sm} = [rectangle, minimum width=0.4cm, minimum height=0.4cm, draw=black, fill=green!30]
        
        % SM grid (simplified)
        \foreach \x in {0,0.5,1,1.5,2,2.5,3,3.5} {
            \foreach \y in {0,0.5,1,1.5,2,2.5,3,3.5} {
                \node[sm] at (\x,\y) {};
            }
        }
        
        \node[below, font=\small] at (1.75,-0.5) {Streaming Multiprocessors (SMs)};
        \node[above, font=\small\bfseries] at (1.75,4.2) {GPU: Miles de núcleos pequeños};
        
        % Comparison arrow
        \draw[ultra thick, <->] (5,1.75) -- (6.5,1.75);
        
        % CPU representation
        \foreach \x in {7.5,8.5} {
            \foreach \y in {1,2.5} {
                \node[rectangle, minimum width=0.8cm, minimum height=0.8cm, draw=black, fill=blue!30] at (\x,\y) {};
            }
        }
        \node[below, font=\small] at (8,-0.5) {Núcleos CPU};
        \node[above, font=\small\bfseries] at (8,4.2) {CPU: Pocos núcleos grandes};
    \end{tikzpicture}
    
    \vspace{0.3cm}
    \begin{columns}
        \column{0.5\textwidth}
        \centering
        \textbf{NVIDIA H100:} 16,896 núcleos CUDA
        
        \column{0.5\textwidth}
        \centering
        \textbf{Intel i9-14900K:} 24 núcleos
    \end{columns}
\end{frame}

\begin{frame}{Jerarquía de GPUs NVIDIA para IA}
    \begin{table}[]
        \centering
        \renewcommand{\arraystretch}{1.3}
        \scriptsize
        \begin{tabular}{p{2.2cm} p{1.8cm} p{1.5cm} p{2.2cm} p{2cm} p{2cm}}
            \toprule
            \textbf{GPU} & \textbf{VRAM} & \textbf{TDP} & \textbf{Tensor TFLOPS} & \textbf{Precio} & \textbf{Uso típico} \\
            \midrule
            RTX 4090 & 24 GB & 450W & 330 (FP16) & \$1,600 & Entusiastas/Dev \\
            A100 & 40/80 GB & 400W & 312 (TF32) & \$10,000+ & Data Centers \\
            H100 SXM & 80 GB & 700W & 989 (TF32) & \$30,000+ & Entrenam. LLMs \\
            H200 & 141 GB & 700W & 989 (TF32) & \$40,000+ & LLMs grandes \\
            B200 & 192 GB & 1000W & 2,250 (FP8) & \$50,000+ & Frontier AI \\
            \bottomrule
        \end{tabular}
    \end{table}
    
    \vspace{0.3cm}
    \begin{alertblock}{La Memoria es Crítica}
        Un modelo de 70B parámetros en FP16 requiere \textbf{140 GB} de VRAM solo para los pesos. La VRAM es frecuentemente el factor limitante.
    \end{alertblock}
\end{frame}

%===========================================================
% SECCIÓN 4: TPU
%===========================================================
\section{TPU: El Chip de Google}

\begin{frame}{Tensor Processing Unit}
    \begin{columns}
        \column{0.55\textwidth}
        \textbf{¿Qué es una TPU?}
        \begin{itemize}
            \item ASIC diseñado por Google (2016)
            \item Optimizado específicamente para TensorFlow
            \item Disponible en Google Cloud
            \item Arquitectura de matriz sistólica
        \end{itemize}
        
        \vspace{0.3cm}
        \textbf{Generaciones:}
        \begin{itemize}
            \item TPU v1: Solo inferencia
            \item TPU v2/v3: Entrenamiento + inferencia
            \item TPU v4: Pods de hasta 4,096 chips
            \item TPU v5e/v5p: Última generación (2023)
        \end{itemize}
        
        \column{0.45\textwidth}
        \centering
        \begin{tikzpicture}
            % Systolic array visualization
            \foreach \x in {0,0.6,1.2,1.8} {
                \foreach \y in {0,0.6,1.2,1.8} {
                    \node[draw, rectangle, minimum size=0.5cm, fill=orange!30] at (\x,\y) {};
                }
            }
            
            % Arrows showing data flow
            \foreach \y in {0,0.6,1.2,1.8} {
                \draw[->, thick, blue] (-0.5,\y) -- (0,\y);
            }
            \foreach \x in {0,0.6,1.2,1.8} {
                \draw[->, thick, red] (\x,2.3) -- (\x,1.8);
            }
            
            \node[below, font=\scriptsize] at (0.9,-0.4) {Matriz Sistólica};
            \node[left, font=\tiny, blue] at (-0.5,0.9) {Datos};
            \node[above, font=\tiny, red] at (0.9,2.3) {Pesos};
        \end{tikzpicture}
        
        \vspace{0.3cm}
        \scriptsize Los datos ``fluyen'' a través de la matriz
    \end{columns}
\end{frame}

\begin{frame}{Matriz Sistólica: El Corazón de la TPU}
    \begin{block}{Concepto}
        Una matriz sistólica es un arreglo de procesadores donde los datos fluyen rítmicamente, como la sangre en el corazón (sístole).
    \end{block}
    
    \centering
    \begin{tikzpicture}[scale=0.8]
        % Grid of PEs
        \foreach \x in {0,1.5,3,4.5} {
            \foreach \y in {0,1.2,2.4} {
                \node[draw, circle, minimum size=0.8cm, fill=orange!20, font=\tiny] (pe\x\y) at (\x,\y) {MAC};
            }
        }
        
        % Horizontal connections
        \foreach \y in {0,1.2,2.4} {
            \draw[->, thick] (-0.8,\y) -- (0,\y);
            \draw[->, thick] (0.4,\y) -- (1.1,\y);
            \draw[->, thick] (1.9,\y) -- (2.6,\y);
            \draw[->, thick] (3.4,\y) -- (4.1,\y);
            \draw[->, thick] (4.9,\y) -- (5.5,\y);
        }
        
        % Vertical connections  
        \foreach \x in {0,1.5,3,4.5} {
            \draw[->, thick] (\x,3.2) -- (\x,2.8);
            \draw[->, thick] (\x,2) -- (\x,1.6);
            \draw[->, thick] (\x,0.8) -- (\x,0.4);
        }
    \end{tikzpicture}
    
    \vspace{0.3cm}
    \textbf{Ventaja:} Alta eficiencia energética y reutilización de datos. Menos accesos a memoria.
\end{frame}

\begin{frame}{TPU vs GPU: Comparación}
    \begin{table}[]
        \centering
        \renewcommand{\arraystretch}{1.4}
        \begin{tabular}{p{4cm} p{4.5cm} p{4.5cm}}
            \toprule
            \textbf{Aspecto} & \textbf{GPU (NVIDIA)} & \textbf{TPU (Google)} \\
            \midrule
            Flexibilidad & Alta (cualquier framework) & Media (mejor con JAX/TF) \\
            Disponibilidad & Compra o cloud & Solo Google Cloud \\
            Programación & CUDA, PyTorch, etc. & JAX, TensorFlow \\
            Costo/hora (cloud) & \$2-30/hr & \$1.5-8/hr \\
            \bottomrule
        \end{tabular}
    \end{table}
    
    \vspace{0.3cm}
    \begin{exampleblock}{Dato Curioso}
        Google usó TPU v4 pods para entrenar PaLM (540B parámetros) y Gemini.
    \end{exampleblock}
\end{frame}

%===========================================================
% SECCIÓN 5: LPU
%===========================================================
\section{LPU: La Nueva Frontera}

\begin{frame}{Groq LPU: Inferencia Ultrarrápida}
    \begin{columns}
        \column{0.55\textwidth}
        \textbf{¿Qué es una LPU?}
        \begin{itemize}
            \item Language Processing Unit (Groq, 2024)
            \item Diseñada exclusivamente para \textbf{inferencia} de LLMs
        \end{itemize}
        
        \vspace{0.3cm}
        \textbf{Innovación clave:}
        \begin{itemize}
            \item Elimina el cuello de botella de memoria
            \item Latencia determinística
            \item Flujo de datos sincronizado por software
        \end{itemize}
        
        \column{0.45\textwidth}
        \centering
        \begin{tikzpicture}
            \node[draw, rectangle, rounded corners, minimum width=3cm, minimum height=2cm, fill=purple!20] (lpu) {};
            \node[above, font=\bfseries] at (lpu.north) {LPU};
            
            \node[draw, rectangle, minimum width=2.5cm, minimum height=0.5cm, fill=yellow!30, font=\scriptsize] at (0,0.5) {SRAM (On-chip)};
            \node[draw, rectangle, minimum width=2.5cm, minimum height=0.5cm, fill=green!30, font=\scriptsize] at (0,-0.2) {Vector Units};
            \node[draw, rectangle, minimum width=2.5cm, minimum height=0.5cm, fill=blue!30, font=\scriptsize] at (0,-0.9) {Matrix Units};
        \end{tikzpicture}
        
        \vspace{0.3cm}
        \scriptsize Sin memoria externa HBM
    \end{columns}
\end{frame}

\begin{frame}{Rendimiento de Groq}
    \begin{block}{Velocidad de Inferencia (tokens/segundo)}
        Groq LPU puede generar \textbf{500+ tokens/segundo} en modelos como LLaMA-2 70B.
    \end{block}
    
    \centering
    \begin{tikzpicture}
        % Bar chart
        \draw[thick, ->] (0,0) -- (8,0) node[right, font=\scriptsize] {tokens/s};
        \draw[thick, ->] (0,0) -- (0,3.5);
        
        % Bars
        \fill[blue!50] (0.5,0) rectangle (1.5,0.6);
        \node[below, font=\tiny] at (1,0) {CPU};
        \node[above, font=\tiny] at (1,0.6) {~10};
        
        \fill[green!50] (2,0) rectangle (3,1.5);
        \node[below, font=\tiny] at (2.5,0) {RTX 4090};
        \node[above, font=\tiny] at (2.5,1.5) {~50};
        
        \fill[orange!50] (3.5,0) rectangle (4.5,2.1);
        \node[below, font=\tiny] at (4,0) {H100};
        \node[above, font=\tiny] at (4,2.1) {~100};
        
        \fill[purple!50] (5,0) rectangle (6,3.2);
        \node[below, font=\tiny] at (5.5,0) {Groq LPU};
        \node[above, font=\tiny] at (5.5,3.2) {\textbf{500+}};
    \end{tikzpicture}
    
    \vspace{0.3cm}
    \begin{alertblock}{Limitación}
        Las LPUs actuales son solo para \textbf{inferencia}, no para entrenamiento.
    \end{alertblock}
\end{frame}

%===========================================================
% SECCIÓN 6: COMPARACIÓN GLOBAL
%===========================================================
\section{Comparación y Casos de Uso}

\begin{frame}{Tabla Comparativa Completa}
    \begin{table}[]
        \centering
        \renewcommand{\arraystretch}{1.2}
        \scriptsize
        \begin{tabular}{p{2.3cm} p{2.5cm} p{2.8cm} p{2.8cm} p{2.8cm}}
            \toprule
            \textbf{Característica} & \textbf{CPU} & \textbf{GPU} & \textbf{TPU} & \textbf{LPU} \\
            \midrule
            Núcleos & 4-128 & 1,000-20,000 & Matriz sistólica & TSP \\
            Paralelismo & Bajo & Muy alto & Muy alto & Alto (streaming) \\
            Memoria & DDR5 & HBM (80-192GB) & HBM & SRAM on-chip \\
            Entrenamiento & Posible (lento) & \textbf{Excelente} & \textbf{Excelente} & No \\
            Inferencia & Modelos pequeños & Buena & Buena & \textbf{Excepcional} \\
            Flexibilidad & \textbf{Máxima} & Alta & Media & Baja (solo LLM) \\
            Disponibilidad & Universal & Alta & Solo GCP & Limitada (Groq) \\
            Costo inicial & Bajo & Medio-Alto & N/A (cloud) & N/A (cloud) \\
            \bottomrule
        \end{tabular}
    \end{table}
\end{frame}

\begin{frame}{¿Cuándo Usar Cada Uno?}
    \begin{columns}[t]
        \column{0.48\textwidth}
        \begin{block}{CPU \faServer}
            \begin{itemize}
                \item Preprocesamiento de datos
                \item Modelos pequeños (sklearn)
                \item Cuando no hay GPU
            \end{itemize}
        \end{block}
        
        \begin{block}{GPU \faMicrochip}
            \begin{itemize}
                \item Entrenamiento de cualquier modelo
                \item PyTorch / cualquier framework
                \item Fine-tuning de LLMs
            \end{itemize}
        \end{block}
        
        \column{0.48\textwidth}
        \begin{block}{TPU \faGoogle}
            \begin{itemize}
                \item Modelos con TensorFlow/JAX
                \item Cuando costo/rendimiento importa
                \item Cargas de trabajo predecibles
            \end{itemize}
        \end{block}
        
        \begin{block}{LPU \faBolt}
            \begin{itemize}
                \item Inferencia de LLMs en producción
                \item Aplicaciones real-time
                \item Chatbots de alta velocidad
                \item Cuando latencia es crítica
            \end{itemize}
        \end{block}
    \end{columns}
\end{frame}

\begin{frame}{Pipeline Típico de IA}
    \centering
    \begin{tikzpicture}[node distance=0.8cm, auto]
        \tikzstyle{stage} = [rectangle, rounded corners, minimum width=2.2cm, minimum height=1cm, text centered, draw=black, font=\scriptsize]
        \tikzstyle{hw} = [rectangle, minimum width=2cm, minimum height=0.5cm, font=\tiny, fill=gray!20]
        \tikzstyle{arrow} = [thick,->,>=stealth]
        
        \node (data) [stage, fill=blue!20] {Recolección\\de Datos};
        \node (prep) [stage, right=0.6cm of data, fill=blue!20] {Preprocesa-\\miento};
        \node (train) [stage, right=0.6cm of prep, fill=green!20] {Entrena-\\miento};
        \node (eval) [stage, right=0.6cm of train, fill=yellow!20] {Evaluación};
        \node (deploy) [stage, right=0.6cm of eval, fill=purple!20] {Deploy /\\Inferencia};
        
        \draw [arrow] (data) -- (prep);
        \draw [arrow] (prep) -- (train);
        \draw [arrow] (train) -- (eval);
        \draw [arrow] (eval) -- (deploy);
        
        % Hardware labels
        \node[hw, below=0.3cm of data] {CPU};
        \node[hw, below=0.3cm of prep] {CPU};
        \node[hw, below=0.3cm of train] {GPU/TPU};
        \node[hw, below=0.3cm of eval] {GPU};
        \node[hw, below=0.3cm of deploy] {GPU/LPU};
    \end{tikzpicture}
    
    \vspace{0.5cm}
    \begin{exampleblock}{Insight}
        No todo el pipeline necesita el mismo hardware. Optimizar cada etapa por separado puede reducir costos significativamente.
    \end{exampleblock}
\end{frame}

%===========================================================
% SECCIÓN 7: FUTURO Y TENDENCIAS
%===========================================================
\section{El Futuro del Hardware para IA}

\begin{frame}{Tendencias Emergentes}
    \begin{columns}[t]
        \column{0.48\textwidth}
        \begin{block}{Chips Especializados}
            \begin{itemize}
                \item AMD MI300X (competencia a NVIDIA)
                \item Intel Gaudi 3
                \item AWS Trainium/Inferentia
                \item Cerebras (wafer-scale)
            \end{itemize}
        \end{block}
        
        \begin{block}{Computación Neuromórfica}
            \begin{itemize}
                \item Intel Loihi
                \item Imita el cerebro biológico
                \item Ultra bajo consumo
                \item Ideal para edge AI
            \end{itemize}
        \end{block}
        
        \column{0.48\textwidth}
        \begin{block}{Computación Cuántica}
            \begin{itemize}
                \item Potencial para ciertos algoritmos ML
                \item Aún en etapa experimental
                \item IBM, Google, IonQ
            \end{itemize}
        \end{block}
        
        \begin{block}{Computación en el Edge}
            \begin{itemize}
                \item NPUs en smartphones
                \item Apple Neural Engine
                \item Qualcomm AI Engine
                \item Google Edge TPU
            \end{itemize}
        \end{block}
    \end{columns}
\end{frame}

\begin{frame}{La Carrera por el Hardware de IA}
    \centering
    \begin{tikzpicture}
        \tikzstyle{company} = [rectangle, rounded corners, minimum width=2cm, minimum height=0.8cm, font=\scriptsize\bfseries, draw=black]
        
        % Leader
        \node[company, fill=green!40] (nvidia) at (0,0) {NVIDIA};
        \node[right=0.3cm of nvidia, font=\scriptsize] {Dominante (80\%+ market share)};
        
        % Challengers
        \node[company, fill=yellow!40] (amd) at (0,-1) {AMD};
        \node[company, fill=yellow!40] (google) at (3,-1) {Google (TPU)};
        
        % Emerging
        \node[company, fill=orange!30] (groq) at (0,-2) {Groq};
        \node[company, fill=orange!30] (cerebras) at (2.5,-2) {Cerebras};
        \node[company, fill=orange!30] (aws) at (5,-2) {AWS};
        
        % Labels
        \node[left=1cm of nvidia, font=\scriptsize\bfseries] {Líder};
        \node[left=1cm of amd, font=\scriptsize\bfseries] {Retadores};
        \node[left=1cm of groq, font=\scriptsize\bfseries] {Emergentes};
        
        % Future
        \node[company, fill=purple!20] (quantum) at (2.5,-3.2) {¿Quantum?};
        \node[left=1cm of quantum, font=\scriptsize\bfseries] {Futuro};
    \end{tikzpicture}
    
    \vspace{0.3cm}
    \begin{alertblock}{Implicación Geopolítica}
        La producción de chips avanzados está concentrada en Taiwán (TSMC). Esto tiene implicaciones estratégicas globales.
    \end{alertblock}
\end{frame}

%===========================================================
% SECCIÓN 8: PRÁCTICA
%===========================================================
\section{Ejercicios Propuestos}

\begin{frame}{Ejercicios}
    \begin{enumerate}
        \setlength\itemsep{0.7em}
        \item \textbf{Investigación:} Comparar las especificaciones de NVIDIA H100 vs AMD MI300X vs Google TPU v5p. Crear una tabla comparativa.
        
        \item \textbf{Cálculo:} Si un modelo tiene 7B parámetros en FP16, ¿cuánta VRAM necesita? ¿Cabe en una RTX 4090?
        
        \item \textbf{Análisis:} Para un chatbot de servicio al cliente que debe responder en $<$200ms, ¿qué hardware recomendarías y por qué?
        
        \item \textbf{Práctica:} Usar Google Colab para comparar el tiempo de entrenamiento de un modelo simple en CPU vs GPU (T4).
        
        \item \textbf{(Bonus)} Investigar qué es un ``TPU Pod'' y cómo se conectan múltiples TPUs.
    \end{enumerate}
\end{frame}

\begin{frame}{Recursos para Profundizar}
    \begin{columns}
        \column{0.48\textwidth}
        \textbf{Documentación Oficial:}
        \begin{itemize}
            \item NVIDIA CUDA Programming Guide
            \item Google Cloud TPU Documentation
            \item Groq Developer Hub
        \end{itemize}
        
        \vspace{0.3cm}
        \textbf{Papers Fundamentales:}
        \begin{itemize}
            \item ``In-Datacenter Performance Analysis of a TPU'' (Google, 2017)
            \item ``Dissecting the NVIDIA Hopper Architecture'' (2022)
        \end{itemize}
        
        \column{0.48\textwidth}
        \textbf{Herramientas Prácticas:}
        \begin{itemize}
            \item Google Colab (GPU/TPU gratis)
            \item Lambda Labs (cloud GPUs)
            \item vast.ai (GPUs económicas)
            \item Groq Cloud (LPU API)
        \end{itemize}
        
        \vspace{0.3cm}
        \textbf{Canales/Blogs:}
        \begin{itemize}
            \item Two Minute Papers
            \item Chip Huyen's Blog
            \item NVIDIA Developer Blog
        \end{itemize}
    \end{columns}
\end{frame}

%--- Diapositiva Final ---
\begin{frame}
    \centering
    \Huge \textbf{¿Preguntas?}
    
    \vspace{0.8cm}
    \large
    \textbf{Próxima clase:} Bases de Datos: SQL y NoSQL
    
    \vspace{0.5cm}
    \normalsize
    Prof. Luis Alfredo Alvarado Rodríguez\\
    \small Escuela de Ciencias Físicas y Matemáticas
    
    \vspace{0.5cm}
    \footnotesize
    \textit{``El hardware define los límites de lo posible en IA.''}
\end{frame}

\end{document}